\documentclass[11pt,a4paper]{article}

\usepackage[polish]{babel}
\usepackage[T1]{fontenc}
\usepackage[utf8]{inputenc}
\usepackage{lmodern}
\usepackage{booktabs}
\usepackage[table]{xcolor}
\usepackage[margin=2.3cm]{geometry}
\usepackage{tcolorbox}
\usepackage{enumitem}
\usepackage{longtable}
\usepackage[hidelinks]{hyperref}
\usepackage{titlesec}
\usepackage{parskip}

\definecolor{granatowy}{HTML}{1B365D}
\definecolor{ciemnyniebieski}{HTML}{16213e}
\definecolor{zolty}{HTML}{fff3cd}
\definecolor{niebieski}{HTML}{e7f3ff}

\titleformat{\section}{\Large\bfseries\color{granatowy}}{\thesection.}{0.5em}{}
\titleformat{\subsection}{\large\bfseries\color{ciemnyniebieski}}{\thesubsection}{0.5em}{}

\tcbuselibrary{skins,breakable}
\newtcolorbox{zastrzezenie}{
	colback=zolty, colframe=zolty, sharp corners, boxrule=0pt,
	left=8pt, right=8pt, top=6pt, bottom=6pt, breakable
}
\newtcolorbox{infobox}{
	colback=niebieski, colframe=niebieski, sharp corners, boxrule=0pt,
	left=8pt, right=8pt, top=6pt, bottom=6pt, breakable
}

\title{\textbf{Przegląd rynku: ładowanie pojazdów elektrycznych}\\
	\large Ceny, transakcje i mnożniki --- Polska i Europa}
\author{Opracowanie na podstawie monitoringu cenników operatorów\\
	oraz bazy transakcji M\&A i rund finansowania}
\date{Stan na sierpień 2026}

\begin{document}
	\maketitle
	\vspace{-0.8cm}
	\noindent\hrulefill
	\vspace{0.6cm}
	
	\section{Podsumowanie kluczowych wniosków}
	
	\begin{infobox}
		\begin{enumerate}[leftmargin=1.3em, itemsep=4pt]
			\item \textbf{Polski rynek rośnie szybko i strukturalnie} --- flota BEV wzrosła
			o 46\% r/r (styczeń--kwiecień 2026), \emph{mimo braku rządowych dopłat}, co
			sugeruje popyt napędzany rynkiem, nie subsydiami.
			\item \textbf{Infrastruktura przesuwa się w stronę DC} --- udział szybkich
			punktów wzrósł z 37\% (koniec 2025) do 49\% (czerwiec 2026); w I półroczu
			2026 operatorzy uruchomili ponad 1000 nowych punktów DC.
			\item \textbf{Ceny odchodzą od prostych stawek AC/DC} --- co najmniej czterech
			z siedmiu monitorowanych operatorów stosuje cenniki dynamiczne lub
			wielopoziomowe (pora dnia, sezon, próg mocy, kanał płatności).
			\item \textbf{Finansowanie przechodzi z venture capital na dług
				instytucjonalny} --- w latach 2025--2026 dominują green loans i facilities
			od banków, nie rundy equity od funduszy VC.
			\item \textbf{Kapitał instytucjonalny wszedł do Europy Środkowej} ---
			finansowanie GreenWay (138 mln EUR) obejmujące Polskę, Słowację i Chorwację
			to sygnał, że region przestał być postrzegany jako zbyt niedojrzały.
			\item \textbf{Wyceny zależą od utylizacji, nie od liczby punktów} ---
			benchmarki branżowe operują na EV/EBITDA (6--9x dla operatorów sieci),
			a nie na mnożnikach per punkt czy per kW.
		\end{enumerate}
	\end{infobox}
	
	\section{Rynek polski --- skala i dynamika}
	
	\subsection{Flota pojazdów}
	
	Na koniec czerwca 2026 r.\ po polskich drogach jeździło \textbf{292\,472}
	samochody osobowe z napędem elektrycznym: 142\,130 w pełni elektrycznych
	(BEV) oraz 150\,342 hybryd plug-in (PHEV). Do tego dochodzi 13\,817
	samochodów dostawczych i ciężarowych oraz 31\,180 motorowerów i motocykli.
	
	\begin{center}
		\begin{tabular}{lrrr}
			\toprule
			\rowcolor{granatowy}
			\textcolor{white}{\textbf{Moment}} & \textcolor{white}{\textbf{BEV osobowe}} &
			\textcolor{white}{\textbf{Punkty ładowania}} & \textcolor{white}{\textbf{Udział DC}} \\
			\midrule
			Koniec 2025 & ok. 120 tys. & 11\,891 & 37\% \\
			Styczeń 2026 & 125\,602 & 12\,311 & --- \\
			Marzec 2026 & 131\,863 & 12\,543 & 48\% \\
			Czerwiec 2026 & 142\,130 & 13\,201 & 49\% \\
			\bottomrule
		\end{tabular}
	\end{center}
	
	\begin{zastrzezenie}
		\textbf{Uwaga metodologiczna:} od stycznia 2026 r.\ PSNM/PZPM zmieniły sposób
		klasyfikacji punktów AC/DC w stacjach wielozłączowych (przyjęto zasadę
		nadrzędności punktu o wyższej mocy). Część skokowego wzrostu udziału DC
		z 37\% na 48\% wynika więc ze \emph{zmiany metodologii}, nie wyłącznie
		z realnej rozbudowy sieci. Porównania udziału DC przed i po styczniu 2026
		r.\ należy traktować ostrożnie.
	\end{zastrzezenie}
	
	\subsection{Wzrost bez wsparcia publicznego}
	
	Najbardziej znaczący sygnał z pierwszego półrocza 2026 r.: mimo
	\textbf{nieprzyznania} 300 mln zł na dopłaty do pojazdów zeroemisyjnych,
	rejestracje BEV w okresie styczeń--kwiecień wzrosły o \textbf{46\%} względem
	analogicznego okresu 2025 r. Po okresie hossy napędzanej programem
	,,NaszEauto'' rynek wrócił do tempa ok.\ 50\% wzrostu r/r --- czyli poziomu
	obserwowanego przed uruchomieniem dopłat.
	
	Dla inwestora oznacza to, że popyt na infrastrukturę ładowania w Polsce
	opiera się na strukturalnym trendzie rynkowym, a nie na jednorazowym
	impulsie regulacyjnym.
	
	\section{Ceny --- co pokazuje monitoring cenników}
	
	\subsection{Rozpiętość stawek}
	
	Monitoring obejmuje siedmiu operatorów działających w Polsce (ORLEN Charge,
	GreenWay, Budimex Mobility, NOXO, Tauron/eTAURON, Eleport, MOYA energia).
	Rozpiętość stawek na rynku jest znaczna --- od ok.\ 1,29 zł/kWh (taryfa nocna
	AC) do ponad 3,30 zł/kWh (DC u najdroższych operatorów).
	
	\subsection{Główny trend: odejście od prostego cennika AC/DC}
	
	Najważniejsza obserwacja z monitoringu nie dotyczy poziomu cen, lecz ich
	\textbf{struktury}. Klasyczny model ,,jedna stawka AC, jedna stawka DC''
	przestaje być standardem:
	
	\begin{itemize}[leftmargin=1.3em]
		\item \textbf{Tauron} --- cennik dynamiczny powiązany z giełdowymi notowaniami
		energii, publikowany z jednodniowym wyprzedzeniem, plus osobna oferta
		abonamentowa z podziałem na strefę letnią i zimową oraz przedziały godzinowe
		(od 0,29 zł/kWh w weekendy do 2,09 zł/kWh zimą na DC).
		\item \textbf{MOYA energia} --- pięć progów mocy DC (60/120/180/400 kW),
		osobna taryfa dzienna i nocna dla AC, dopłata 0,30 zł/kWh przy płatności
		terminalem zamiast aplikacją.
		\item \textbf{Eleport} --- trzy progi mocy DC z rosnącą stawką.
		\item \textbf{GreenWay} --- trzy plany abonamentowe (0 / 29,99 / 79,99 zł
		miesięcznie) z malejącą stawką za kWh.
	\end{itemize}
	
	\begin{infobox}
		\textbf{Konsekwencja praktyczna:} porównywanie operatorów pojedynczą liczbą
		staje się coraz mniej miarodajne. Realna stawka zależy od pory dnia, mocy
		konkretnej stacji, kanału płatności i posiadania abonamentu. Ma to też
		bezpośrednie przełożenie na modelowanie przychodów z nowej lokalizacji ---
		przyjęcie jednej, uśrednionej ceny za kWh coraz gorzej odzwierciedla
		faktyczną strukturę rynku.
	\end{infobox}
	
	\subsection{Rola rabatów i programów lojalnościowych}
	
	Operatorzy powiązani z sieciami stacji paliw (ORLEN, MOYA) integrują
	ładowanie z istniejącymi programami lojalnościowymi, umożliwiając obniżenie
	kosztu sesji nawet o 50\% przez punkty. To zaciera obraz cenowy jeszcze
	bardziej --- efektywna cena płacona przez stałego klienta może istotnie
	odbiegać od publikowanej stawki.
	
	\section{Transakcje --- gdzie płynie kapitał}
	
	\subsection{Skala aktywności}
	
	Baza obejmuje \textbf{27 pozycji}: 18 rund finansowania i 9 transakcji
	M\&A/fuzji, z lat 2018--2026, w Niemczech, Holandii, Francji, Hiszpanii,
	Norwegii, Portugalii, Wielkiej Brytanii i Polsce. Same rundy finansowania
	z lat 2022--2026 opiewają na kwoty rzędu \textbf{4,7 mld EUR}.
	
	\subsection{Trend 1: od equity do długu}
	
	Rundy z lat 2021--2023 miały charakter kapitałowy (IONITY: 700 mln EUR
	equity od GIP/BlackRock; Milence: 500 mln EUR od Daimler Truck, TRATON
	i Volvo; Driveco: 250 mln EUR). Od 2024--2025 dominują natomiast
	\textbf{instrumenty dłużne}: green loans, senior debt facilities,
	convertible loans.
	
	Rekordowa transakcja --- \textbf{600 mln EUR green loan dla IONITY} (maj
	2025), udzielony przez konsorcjum dziewięciu banków --- jest największym
	finansowaniem w historii europejskiej branży ładowania.
	
	\begin{infobox}
		\textbf{Co to oznacza:} banki i kredytodawcy instytucjonalni zaczęli
		traktować infrastrukturę ładowania jako \emph{klasę aktywów
			infrastrukturalnych}, a nie inwestycję wysokiego ryzyka. Dług wymaga
		jednak przewidywalnych przepływów pieniężnych, co przesuwa uwagę
		inwestorów z tempa rozbudowy sieci na \textbf{utylizację i jakość
			lokalizacji}.
	\end{infobox}
	
	\subsection{Trend 2: kapitał wchodzi do Europy Środkowej}
	
	Finansowanie \textbf{GreenWay} (138 mln EUR green debt, 2025--2026),
	obejmujące Polskę, Słowację i Chorwację, jest sygnałem odmiennym od
	pozostałych transakcji w bazie --- wszystkie pozostałe dotyczą Europy
	Zachodniej lub Skandynawii. Oznacza to, że region CEE przestał być
	postrzegany przez kredytodawców jako zbyt niedojrzały do finansowania
	dłużnego.
	
	Dodatkowo \textbf{Powerdot} (Portugalia) wprost wymienia Polskę wśród
	swoich rynków działania, co potwierdza zainteresowanie zagranicznych
	operatorów polskim rynkiem.
	
	\subsection{Trend 3: konsolidacja przez koncerny paliwowe i energetyczne}
	
	Segment M\&A jest zdominowany przez przejęcia dokonywane przez duże
	podmioty spoza branży:
	
	\begin{center}
		\begin{tabular}{llr}
			\toprule
			\rowcolor{granatowy}
			\textcolor{white}{\textbf{Nabywca}} & \textcolor{white}{\textbf{Cel}} & \textcolor{white}{\textbf{Kwota}} \\
			\midrule
			BP & Chargemaster (UK) & ok. 147 mln EUR \\
			EDF + Legal \& General & Pod Point (53\%, UK) & ok. 125 mln EUR \\
			Shell & ubitricity (DE/UK) & nieujawniona \\
			Shell & Volta (USA) & ok. 155 mln EUR \\
			TotalEnergies & Blue Charge (Singapur) & nieujawniona \\
			Mer (Statkraft) & eeMobility (DE) & nieujawniona \\
			Wallbox & ABL (DE) & 15 mln EUR \\
			\bottomrule
		\end{tabular}
	\end{center}
	
	Motywacja jest dwojaka: wejście na nowy rynek geograficzny (TotalEnergies
	w Singapurze, Mer w Niemczech) lub uzupełnienie oferty detalicznej
	o ładowanie (BP, Shell, EDF).
	
	\subsection{Trend 4: korekta wycen po szczycie SPAC}
	
	Przejęcie \textbf{Volta przez Shell} (2023, ok.\ 170 mln USD) nastąpiło przy
	wycenie niższej o ok.\ \textbf{95\%} względem szczytowej wyceny z debiutu
	giełdowego przez SPAC w 2021 r. Notowane spółki branżowe (ChargePoint,
	Blink, EVgo) handlowane są znacznie poniżej wycen z IPO. Rynek przeszedł
	z fazy spekulacyjnej do fazy oceny fundamentalnej.
	
	\section{Mnożniki wyceny}
	
	\subsection{Czego branża nie publikuje}
	
	Przeszukano dostępne źródła pod kątem mnożników \textbf{EV/kW}
	i \textbf{EV/punkt ładowania}. Takie wskaźniki \textbf{nie są publikowane}
	jako zagregowane benchmarki branżowe, a dla poszczególnych transakcji
	źródła zwykle ujawniają kwotę \emph{albo} wielkość sieci, rzadko oba
	naraz i z tego samego momentu.
	
	Ma to logiczne uzasadnienie: kluczowym czynnikiem wyceny jest
	\textbf{utylizacja}, a nie sama liczba czy moc zainstalowanych punktów.
	Dwie sieci o identycznej liczbie ładowarek mogą mieć skrajnie różną
	wartość, jeśli różnią się wykorzystaniem.
	
	\subsection{Co branża publikuje: EV/EBITDA}
	
	\begin{center}
		\begin{tabular}{lc}
			\toprule
			\rowcolor{granatowy}
			\textcolor{white}{\textbf{Profil biznesowy}} & \textcolor{white}{\textbf{Mnożnik EV/EBITDA}} \\
			\midrule
			Wykonawcy instalacji & 3,0--5,0x \\
			Instalacja + serwis (regionalnie) & 4,0--6,0x \\
			Operatorzy sieci (CPO) z istotną utylizacją & 6,0--9,0x \\
			Sieci hubów DC & 7,0--10,0x \\
			Platformy premium o dużej skali & 8,0--12,0x+ \\
			\midrule
			Rynek szwajcarski --- mnożniki transakcyjne & 6,5--9,5x \\
			Rynek szwajcarski --- zakres statystyczny & 5,0--7,0x \\
			\bottomrule
		\end{tabular}
	\end{center}
	
	Dwa niezależne źródła (CT Acquisitions, Valindex) podają zbieżne zakresy dla
	operatorów sieci, co wzmacnia wiarygodność tych benchmarków. Różnica między
	mnożnikami transakcyjnymi (6,5--9,5x) a zakresem statystycznym (5,0--7,0x)
	w danych szwajcarskich wskazuje na \textbf{premię płaconą w realnych
		transakcjach} względem wyceny księgowej.
	
	\begin{zastrzezenie}
		\textbf{Ograniczenie mnożników EBITDA w tej branży:} znaczna część
		operatorów --- w tym notowani na giełdzie (ChargePoint: EBITDA $-$75 mln USD)
		--- nie generuje dodatniej EBITDA. Dla takich podmiotów mnożnik EV/EBITDA
		jest nieprzydatny, a wyceny opierają się na przychodach, wartości aktywów
		lub zdyskontowanych przepływach. To dodatkowo tłumaczy, dlaczego branża nie
		wypracowała jednego, powszechnie stosowanego wskaźnika wyceny.
	\end{zastrzezenie}
	
	\section{Wnioski dla projektu}
	
	\begin{enumerate}[leftmargin=1.3em, itemsep=5pt]
		\item \textbf{Modelowanie przychodów wymaga uwzględnienia struktury taryf.}
		Przyjęcie jednej, uśrednionej stawki za kWh coraz gorzej odzwierciedla rynek,
		na którym dominują taryfy godzinowe, sezonowe i zależne od mocy.
		
		\item \textbf{Utylizacja jest kluczową zmienną wyceny.} Zarówno benchmarki
		wyceny, jak i warunki stawiane przez kredytodawców wskazują, że to
		wykorzystanie stacji --- nie ich liczba --- decyduje o wartości aktywa.
		
		\item \textbf{Finansowanie dłużne jest realnie dostępne dla dojrzałych
			platform.} Skala transakcji z lat 2025--2026 pokazuje, że operator
		z udokumentowaną utylizacją i jakością lokalizacji może liczyć na
		finansowanie bankowe, nie tylko kapitał wysokiego ryzyka.
		
		\item \textbf{Polska jest w polu widzenia kapitału instytucjonalnego.}
		Finansowanie GreenWay oraz obecność Powerdot potwierdzają, że rynek CEE
		przeszedł próg wiarygodności dla inwestorów zagranicznych.
		
		\item \textbf{Wyceny wróciły do poziomów fundamentalnych.} Korekta po
		szczycie SPAC (Volta: $-$95\% względem debiutu) oznacza, że obecne mnożniki
		są bardziej realistyczną podstawą do rozmów transakcyjnych niż dane
		z lat 2021--2022.
	\end{enumerate}
	
	\vspace{0.5cm}
	\noindent\hrulefill
	
	\vspace{0.3cm}
	\small\textit{Źródła danych: pliki \texttt{Zestawienie\_Cennikow\_EV.xlsx}
		(monitoring cenników, 7 operatorów) oraz \texttt{baza\_transakcji\_ev.xlsx}
		(27 transakcji + benchmarki wyceny). Każda pozycja w obu plikach zawiera
		link źródłowy i datę dostępu. Dane rynkowe: Licznik Elektromobilności
		PSNM/PZPM (czerwiec 2026).}
	
\end{document}